\begin{table}[H]
    \centering
    \caption{Matriks Translasi Persentil Standar ke Skor Komposit Lintas Skenario Korelasi Antarsubtes ($M = 500$)}
    \label{tab:ranking_komposit}
    \begin{spacing}{1.0}
    \resizebox{\textwidth}{!}{%
    \begin{tabular}{lccccccccc}
        \toprule
        Titik \textit{Benchmark} & Persentil (\%) & Peringkat Nasional & Raw Sample & r = 0.350 & r = 0.400 & Thorndike Case III & r = 0.450 & r = 0.500 & r = 0.550 \\
        \midrule
        M - 3 SD   & 99,87 & 845.375 & 319,63 & 309,83 & 300,45 & 296,85 & 291,49 & 282,90 & 274,63 \\
        M - 2 SD   & 97,72 & 827.260 & 379,75 & 373,22 & 366,97 & 364,57 & 360,99 & 355,26 & 349,75 \\
        M - 1 SD   & 84,13 & 712.213 & 439,88 & 436,61 & 433,48 & 432,28 & 430,50 & 427,63 & 424,88 \\
        M          & 50,00 & 423.259 & 500,00 & 500,00 & 500,00 & 500,00 & 500,00 & 500,00 & 500,00 \\
        M + 1 SD   & 15,87 & 134.305 & 560,12 & 563,39 & 566,52 & 567,72 & 569,50 & 572,37 & 575,12 \\
        M + 2 SD   & 2,28 & 19.258 & 620,25 & 626,78 & 633,03 & 635,43 & 639,01 & 644,74 & 650,25 \\
        M + 3 SD   & 0,13 & 1.143 & 680,37 & 690,17 & 699,55 & 703,15 & 708,51 & 717,10 & 725,37 \\
        \bottomrule
    \end{tabular}%
    }
    \end{spacing}
    \vspace{2mm}
    \begin{flushleft}
    \footnotesize\textit{Catatan.} Persentil (\textit{top \%}) identik lintas skenario pada baris yang sama secara desain. Hanya skor komposit yang berbeda antarskenario karena perbedaan $SD$ skor komposit teoretis di bawah tiap asumsi korelasi. Thorndike Case III adalah estimasi terkoreksi \textit{restriction of range} (lihat Metode), sedangkan skenario $r$ lain adalah titik pembanding hipotetis dengan asumsi korelasi rerata seragam antarpasangan subtes. Peringkat nasional dihitung dari jumlah peserta hadir, $N_{\text{nasional}} = 846.518$.
    \end{flushleft}
\end{table}
